% Options for packages loaded elsewhere
\PassOptionsToPackage{unicode}{hyperref}
\PassOptionsToPackage{hyphens}{url}
\documentclass[
  11pt,
  a4paper]{article}
\usepackage{xcolor}
\usepackage{amsmath,amssymb}
\setcounter{secnumdepth}{5}
\usepackage{iftex}
\ifPDFTeX
  \usepackage[T1]{fontenc}
  \usepackage[utf8]{inputenc}
  \usepackage{textcomp} % provide euro and other symbols
\else % if luatex or xetex
  \usepackage{unicode-math} % this also loads fontspec
  \defaultfontfeatures{Scale=MatchLowercase}
  \defaultfontfeatures[\rmfamily]{Ligatures=TeX,Scale=1}
\fi
\usepackage{lmodern}
\ifPDFTeX\else
  % xetex/luatex font selection
  \setmainfont[]{Latin Modern Roman}
\fi
% Use upquote if available, for straight quotes in verbatim environments
\IfFileExists{upquote.sty}{\usepackage{upquote}}{}
\IfFileExists{microtype.sty}{% use microtype if available
  \usepackage[]{microtype}
  \UseMicrotypeSet[protrusion]{basicmath} % disable protrusion for tt fonts
}{}
\makeatletter
\@ifundefined{KOMAClassName}{% if non-KOMA class
  \IfFileExists{parskip.sty}{%
    \usepackage{parskip}
  }{% else
    \setlength{\parindent}{0pt}
    \setlength{\parskip}{6pt plus 2pt minus 1pt}}
}{% if KOMA class
  \KOMAoptions{parskip=half}}
\makeatother
\usepackage{longtable,booktabs,array}
\usepackage{calc} % for calculating minipage widths
% Correct order of tables after \paragraph or \subparagraph
\usepackage{etoolbox}
\makeatletter
\patchcmd\longtable{\par}{\if@noskipsec\mbox{}\fi\par}{}{}
\makeatother
% Allow footnotes in longtable head/foot
\IfFileExists{footnotehyper.sty}{\usepackage{footnotehyper}}{\usepackage{footnote}}
\makesavenoteenv{longtable}
\setlength{\emergencystretch}{3em} % prevent overfull lines
\providecommand{\tightlist}{%
  \setlength{\itemsep}{0pt}\setlength{\parskip}{0pt}}

\usepackage{amsmath,amssymb,amsthm}
\usepackage{booktabs}
\usepackage{graphicx}
\usepackage[margin=1in]{geometry}
\usepackage{hyperref}
\usepackage{xcolor}
\usepackage{unicode-math}
\hypersetup{colorlinks=true,linkcolor=blue!60!black,citecolor=green!50!black,urlcolor=blue!60!black}
\usepackage{fancyhdr}
\pagestyle{fancy}
\fancyhf{}
\fancyhead[L]{\small PACSeries Paper \thepapernumber}
\fancyhead[R]{\small Dawn Field Institute}
\fancyfoot[C]{\thepage}
\renewcommand{\headrulewidth}{0.4pt}

% Paper number counter
\newcounter{papernumber}

\setcounter{papernumber}{4}
\usepackage{bookmark}
\IfFileExists{xurl.sty}{\usepackage{xurl}}{} % add URL line breaks if available
\urlstyle{same}
\hypersetup{
  hidelinks,
  pdfcreator={LaTeX via pandoc}}

\title{Standard Model Parameters from Fibonacci Arithmetic}
\author{Peter Groom}
\date{February 2026}

\begin{document}
\maketitle

\renewcommand*\contentsname{Contents}
{
\setcounter{tocdepth}{3}
\tableofcontents
}
\section{Standard Model Parameters from Fibonacci
Arithmetic}\label{standard-model-parameters-from-fibonacci-arithmetic}

\textbf{Peter Groom, Dawn Field Institute}\\
\textbf{PACSeries Paper 4}\\
\textbf{Date}: February 2026\\
\textbf{Version}: 2.1

\begin{center}\rule{0.5\linewidth}{0.5pt}\end{center}

\subsection{Abstract}\label{abstract}

The Standard Model of particle physics contains approximately 25 free
parameters --- coupling constants, mixing angles, and mass ratios ---
whose values are measured but not derived from first principles. This
paper shows that a significant subset of these parameters can be
expressed as closed-form ratios of Fibonacci numbers, with precisions
ranging from 0.5 ppm (Koide formula) to 1.7\% (strong coupling
constant).

The expressions are not fitted. They follow from a single structural
constraint: the PAC recursion \(\Psi(k) = \Psi(k+1) + \Psi(k+2)\), whose
unique stable solution \(\Psi(k) = \varphi^{-k}\) selects the golden
ratio algebraically. The Fibonacci numbers are its integer projections.
The gauge group structure of the Standard Model ---
\(U(1) \times SU(2) \times SU(3)\) --- is the unique combination whose
adjoint dimensions (\(1, 3, 8\)) are Fibonacci numbers, closing at
\(F_7 = 13 = 1 + 3 + 8 + 1\).

Individual Fibonacci matches are not significant (any small integer has
\$\sim\$16\% probability of being Fibonacci). What is significant is the
joint constraint: a single recursion, evaluated at specific hierarchy
depths, simultaneously reproduces gauge couplings to 5.7 ppm, mixing
angles to 0.19\%, mass ratios to 5 ppm, and the Casimir regularisation
factor exactly --- with a combined probability against chance of
\(p < 10^{-5}\). The weak mixing angle prediction
\(\sin^2\theta_W = 3/13\) is excluded at \(M_Z\) (\$\sim\(11\)\sigma\$)
but matches the running value at \(Q = 82.78\) GeV \(\approx M_W\), with
\(M_W/M_Z\) predicted at 0.03\% error (§4.4).

We also derive the She-Lévêque turbulence intermittency constant
\(k = d \times F_{d+1}\) from first principles, connecting particle
physics constants to fluid dynamics through the same Fibonacci
structure. A falsifiable prediction is offered: a Z' boson at
\(395 \pm 20\) GeV with coupling \(g_{Z'}/g_Z = 1/13\).

\textbf{Keywords}: Standard Model, Fibonacci numbers, golden ratio,
coupling constants, mass ratios, PAC conservation, Koide formula, fine
structure constant, Dawn Field Theory

\begin{center}\rule{0.5\linewidth}{0.5pt}\end{center}

\subsection{§1. The Problem of Free
Parameters}\label{the-problem-of-free-parameters}

The Standard Model of particle physics is one of the most precisely
tested theories in science. It predicts the anomalous magnetic moment of
the electron to 12 significant figures. It predicted the existence of
the W and Z bosons, the top quark, and the Higgs boson before they were
discovered.

It does not explain why the fine structure constant \(\alpha\) equals
\(1/137.036\) rather than some other number. It does not explain why
there are three generations of fermions rather than two or seven. It
does not explain why the electron is 206.768 times lighter than the
muon. These values are measured with extraordinary precision but enter
the theory as inputs, not outputs.

This paper asks whether a single structural constraint --- conservation
under hierarchical branching --- can reproduce these values as necessary
consequences.

\begin{center}\rule{0.5\linewidth}{0.5pt}\end{center}

\subsection{§2. The Constraint}\label{the-constraint}

Paper 1 in this series established that information erasure into
multi-mode environments creates correlational structure. When the
environment has cascade topology (sequential mode coupling), the PAC
conservation constraint requires:

\[\Psi(k) = \Psi(k+1) + \Psi(k+2)\]

This is the Fibonacci recursion. Its unique stable solution is
\(\Psi(k) = \varphi^{-k}\), where \(\varphi = (1 + \sqrt{5})/2\) is the
golden ratio. This is not a choice. It is the only bounded solution to
the characteristic equation \(x^2 = x + 1\).

The Fibonacci numbers
\(F_n = \{1, 1, 2, 3, 5, 8, 13, 21, 34, 55, 89, 144, \ldots\}\) are the
integer projections of this recursion. They provide a complete basis for
integer representation (Zeckendorf's theorem: every positive integer has
a unique decomposition into non-consecutive Fibonacci numbers).

Paper 2 established that the golden ratio identity
\(\varphi^2 = \varphi + 1\) holds to \(< 10^{-14}\) across all numerical
bases, confirming that the Fibonacci structure is mathematical, not
representational.

\subsubsection{§2.1 The Lagrangian}\label{the-lagrangian}

The PAC constraint expressed as a Lagrangian density:

\[\mathcal{L} = \frac{1}{2}\sum_k \left(\frac{\partial\Psi_k}{\partial t}\right)^2 - V(\Psi) + \lambda \sum_k [\Psi_k - \Psi_{k+1} - \Psi_{k+2}]^2\]

This has a golden scaling symmetry:
\(\Psi(k) \to \varphi^{-1} \cdot \Psi(k-1)\). By Noether's theorem, this
continuous symmetry yields a conserved charge:

\[Q_{\text{PAC}} = -\frac{1}{\varphi} \int \Psi \cdot \frac{\partial\Psi}{\partial t} \, dx\]

The theory predicts that physical coupling constants are ratios of this
conserved charge evaluated at different hierarchy levels.

\begin{center}\rule{0.5\linewidth}{0.5pt}\end{center}

\subsection{§3. The Gauge Group
Selection}\label{the-gauge-group-selection}

\subsubsection{\texorpdfstring{§3.1 Fibonacci Filter on
\(SU(N)\)}{§3.1 Fibonacci Filter on SU(N)}}\label{fibonacci-filter-on-sun}

A gauge group \(SU(N)\) has \(N^2 - 1\) generators (the dimension of its
adjoint representation). The PAC constraint requires this number to be a
Fibonacci number:

\begin{longtable}[]{@{}lll@{}}
\toprule\noalign{}
Group & \(N^2 - 1\) & Fibonacci? \\
\midrule\noalign{}
\endhead
\bottomrule\noalign{}
\endlastfoot
\(U(1)\) & 1 & \(F_1 = F_2 =\) \textbf{Yes} \\
\(SU(2)\) & 3 & \(F_4 =\) \textbf{Yes} \\
\(SU(3)\) & 8 & \(F_6 =\) \textbf{Yes} \\
\(SU(4)\) & 15 & \textbf{No} \\
\(SU(5)\) & 24 & \textbf{No} \\
\end{longtable}

\(SU(2)\) and \(SU(3)\) are the only non-abelian gauge groups whose
adjoint representations have Fibonacci dimension. This is a constraint,
not a selection --- PAC conservation cannot maintain coherence across a
non-Fibonacci number of coupled modes.

\textbf{Script}: \texttt{exp\_19\_su4\_forbidden.py} --- confirms no
\(SU(N)\) with \(N > 3\) has a Fibonacci adjoint dimension for any
\(N \leq 100\).

\subsubsection{\texorpdfstring{§3.2 PAC Closure at
\(F_7 = 13\)}{§3.2 PAC Closure at F\_7 = 13}}\label{pac-closure-at-f_7-13}

The total gauge content of the Standard Model:

\[U(1): 1 \quad + \quad SU(2): 3 \quad + \quad SU(3): 8 \quad + \quad \text{Higgs}: 1 \quad = \quad 13 = F_7\]

This is PAC closure: the total number of gauge and scalar degrees of
freedom equals a Fibonacci number. The Higgs is counted as 1 physical
degree of freedom --- the radial mode that survives after electroweak
symmetry breaking, where the three Goldstone components are absorbed by
the \(W^\pm\) and \(Z\) bosons. Before EWSB, the Higgs doublet has 4
real components; after EWSB, 1 physical scalar remains. The PAC count
applies to the broken phase --- the phase in which all measured
couplings are defined.

Whether this counting is numerology or structure depends on whether
\(F_7\) predicts anything. It does. The fine structure constant, the
weak mixing angle, the Cabibbo angle, and the neutrino mixing angles all
involve \(F_7 = 13\) in their denominators.

\begin{center}\rule{0.5\linewidth}{0.5pt}\end{center}

\subsection{§4. Gauge Coupling
Constants}\label{gauge-coupling-constants}

\subsubsection{§4.1 Fine Structure
Constant}\label{fine-structure-constant}

The fine structure constant \(\alpha\) measures the strength of
electromagnetic coupling. Its value is known to 0.15 ppb:

\[\alpha^{-1} = 137.035999177(21) \quad \text{(CODATA 2022 [5])}\]

The PAC formula:

\[\alpha = \frac{F_3}{F_4 \cdot \varphi \cdot F_{10}} \left(1 - \frac{F_{10}}{4\pi F_7^2}\right) = \frac{2}{3 \cdot \varphi \cdot 55} \left(1 - \frac{55}{4\pi \cdot 169}\right)\]

\textbf{Computed}: \(\alpha_{\text{PAC}} = 0.00729731\)\\
\textbf{Measured}: \(\alpha_{\text{exp}} = 0.0072973526\)\\
\textbf{Error}: \(5.7 \text{ ppm}\)

\subsubsection{§4.2 Structure of the
Formula}\label{structure-of-the-formula}

The formula has two factors:

\begin{enumerate}
\def\labelenumi{\arabic{enumi}.}
\item
  \textbf{Base rate}:
  \(\frac{F_3}{F_4 \cdot \varphi \cdot F_{10}} = \frac{2}{3\varphi \cdot 55}\).
  This is the coupling strength for one photon propagating through a
  Fibonacci hierarchy of depth \(F_{10} = 55\).
\item
  \textbf{Correction}:
  \(\left(1 - \frac{F_{10}}{4\pi F_7^2}\right) = \left(1 - \frac{55}{4\pi \cdot 169}\right) \approx 0.97393\).
  This accounts for the gauge closure constraint at \(F_7 = 13\).
\end{enumerate}

The interpretation from Paper 1: \(\alpha\) measures the Landauer
structure cost of electromagnetic interaction --- the rate at which
correlational structure \(\xi\) accumulates through 55 hierarchy levels,
corrected by the total gauge content.

\subsubsection{§4.3 Uniqueness}\label{uniqueness}

All Fibonacci pairs \((F_m, F_n)\) for \(m, n \in [3, 20]\) were tested
in the formula. The pair \((10, 7)\) is \(2{,}870\times\) better than
the next best match. Given \(F_{10} = 55\), the formula predicts the
other index must satisfy \(F_n = 13.0014\); the nearest Fibonacci number
\(F_7 = 13\) is 0.001 away.

The continuous Fibonacci extension (\(\varphi^n/\sqrt{5}\)) gives
\emph{worse} results (137 ppm vs 5.7 ppm). Discrete integers are
structurally significant, consistent with the discrete level symmetry of
the PAC Lagrangian.

\textbf{Script}: \texttt{exp\_12\_alpha\_formula.py},
\texttt{exp\_13\_alpha\_falsification.py}

\subsubsection{§4.4 Weak Mixing Angle}\label{weak-mixing-angle}

The Weinberg angle \(\theta_W\) parametrises the electroweak mixing:

\[\sin^2\theta_W = \frac{F_4}{F_7} = \frac{3}{13} = 0.230769\ldots\]

\textbf{Measured}: \(\sin^2\theta_W = 0.23121 \pm 0.00004\) (on-shell,
\(M_Z\)) {[}1{]}\\
\textbf{Error}: \(0.19\%\)

The ratio \(F_4/F_7 = 3/13\) directly gives the weak mixing as the
fraction of \(SU(2)\) generators within the total gauge content. This is
a structural prediction: the mixing is determined by gauge group
dimensions, not by a free parameter.

\textbf{Energy-scale resolution (milestone3/exp\_08).}
\(\sin^2\theta_W\) runs with energy. The on-shell value at
\(M_Z \approx 91.2\) GeV is \(0.23121 \pm 0.00004\), differing from
\(3/13 = 0.23077\) by \$\sim\(11\)\sigma\$. However, the PAC ratio need
not hold at \(M_Z\). A one-loop running calculation finds that
\(\sin^2\theta_W(Q)\) passes through exactly \(3/13\) at:

\[Q = 82.78 \text{ GeV}\]

This is within 3\% of \(M_W = 80.37\) GeV (\(Q/M_W = 1.030\)). The
proximity to \(M_W\) may be physically significant: the W boson mediates
flavour-changing charged-current transitions --- the only Standard Model
process that converts between quark and lepton generations. In PAC
terms, flavour change corresponds to actualization (potential states
becoming actual states), and the Fibonacci mixing \(F_4/F_7\) would then
achieve its exact value at the energy where this actualization mechanism
operates. Whether this interpretation survives beyond the one-loop
approximation used here is an open question.

The tree-level mass ratio prediction follows directly:

\[\frac{M_W}{M_Z} = \cos\theta_W = \sqrt{1 - \frac{3}{13}} = \sqrt{\frac{10}{13}} = 0.8771\]

compared to the on-shell value \(\cos\theta_W = 0.8768\), an error of
\textbf{0.03\%}.

Among 24 Fibonacci ratios \(F_i/F_j\) tested, \(F_4/F_7 = 3/13\) is rank
\#1 for proximity to the measured \(\sin^2\theta_W\).

\textbf{Sensitivity}: \(\pm 3\sigma\) in \(\sin^2\theta_W\) maps to
\(Q \in [81.2, 84.4]\) GeV, bracketing \(M_W\) from above.

\textbf{Scripts}: \texttt{exp\_18\_weinberg\_angle.py},
\texttt{milestone3/scripts/exp\_08\_weinberg\_running.py}

\subsubsection{§4.5 Strong Coupling
Constant}\label{strong-coupling-constant}

\[\alpha_s(M_Z) = \frac{F_4}{2\varphi \cdot F_6} = \frac{3}{2\varphi \cdot 8} = 0.1159\]

\textbf{Measured}: \(\alpha_s(M_Z) = 0.1179 \pm 0.0010\) {[}1{]}\\
\textbf{Error}: \(1.71\%\)

This is the weakest of the three gauge coupling predictions. The formula
gives the strong coupling as the \(SU(2)\) adjoint dimension divided by
twice the golden ratio times the \(SU(3)\) adjoint dimension. The 1.71\%
error is within the measurement uncertainty range but is substantially
less precise than the \(\alpha\) result.

A structural asymmetry deserves comment. The \(\alpha\) formula (§4.1)
includes a correction factor \((1 - F_{10}/4\pi F_7^2)\), while the
\(\alpha_s\) formula is a bare ratio. If all gauge couplings emerge from
the same recursion at different depths, their formulas should share a
template. One interpretation: electromagnetism couples to the full gauge
content (all 13 modes contribute vacuum polarisation loops), requiring
the \(F_7^2\) correction, while the strong coupling operates within a
single gauge sector (\(SU(3)\)) and sees no cross-sector correction.
This is consistent with asymptotic freedom --- at high energy,
\(\alpha_s\) is determined by \(SU(3)\) self-coupling alone. Whether
this interpretation survives a proper renormalisation group analysis is
an open question.

\subsubsection{§4.6 Summary of Gauge
Couplings}\label{summary-of-gauge-couplings}

\begin{longtable}[]{@{}
  >{\raggedright\arraybackslash}p{(\linewidth - 8\tabcolsep) * \real{0.1961}}
  >{\raggedright\arraybackslash}p{(\linewidth - 8\tabcolsep) * \real{0.2549}}
  >{\raggedright\arraybackslash}p{(\linewidth - 8\tabcolsep) * \real{0.2157}}
  >{\raggedright\arraybackslash}p{(\linewidth - 8\tabcolsep) * \real{0.1961}}
  >{\raggedright\arraybackslash}p{(\linewidth - 8\tabcolsep) * \real{0.1373}}@{}}
\toprule\noalign{}
\begin{minipage}[b]{\linewidth}\raggedright
Coupling
\end{minipage} & \begin{minipage}[b]{\linewidth}\raggedright
PAC Formula
\end{minipage} & \begin{minipage}[b]{\linewidth}\raggedright
PAC Value
\end{minipage} & \begin{minipage}[b]{\linewidth}\raggedright
Measured
\end{minipage} & \begin{minipage}[b]{\linewidth}\raggedright
Error
\end{minipage} \\
\midrule\noalign{}
\endhead
\bottomrule\noalign{}
\endlastfoot
\(\alpha\) (EM) &
\(\frac{F_3}{F_4 \cdot \varphi \cdot F_{10}}\left(1 - \frac{F_{10}}{4\pi F_7^2}\right)\)
& 0.0072973 & 0.0072974 & \textbf{5.7 ppm} \\
\(\sin^2\theta_W\) (Weak) & \(F_4/F_7 = 3/13\) & 0.23077 & 0.23121 &
\textbf{0.19\%} \\
\(\alpha_s\) (Strong) & \(F_4/(2\varphi F_6)\) & 0.1159 & 0.1179 &
\textbf{1.71\%} \\
\end{longtable}

The precision ordering (\(\alpha > \sin^2\theta_W > \alpha_s\))
corresponds to the algebraic/transcendental character of the formulas:
the algebraic ratios (\(3/13\)) are exact; the transcendental
corrections (\(\varphi\), \(\pi\)) introduce residual error. Paper 2
discusses the discrete-to-continuous gap as a structural feature, not a
limitation.

\begin{center}\rule{0.5\linewidth}{0.5pt}\end{center}

\subsection{§5. The Koide Formula}\label{the-koide-formula}

The Koide formula is an empirical relation among the charged lepton
masses:

\[Q = \frac{m_e + m_\mu + m_\tau}{(\sqrt{m_e} + \sqrt{m_\mu} + \sqrt{m_\tau})^2} = 0.666661 \pm 0.000007\]

This is \(2/3\) to within 0.5 ppm. The Standard Model does not explain
why {[}2{]}.

In Fibonacci arithmetic:

\[Q = \frac{F_3}{F_3 + F_2} = \frac{2}{2 + 1} = \frac{2}{3}\]

This is not an approximation. It is an exact Fibonacci identity. But the
nontrivial claim is not that \(F_3/(F_3 + F_2) = 2/3\) --- that is
arithmetic. The nontrivial claim is that the Koide ratio maps to this
specific recursion level.

Why \((F_3, F_2)\) and not \((F_4, F_3) = 3/4\) or \((F_5, F_4) = 5/8\)?
Because charged leptons are the shallowest fermionic sector in the PAC
hierarchy. They carry the lowest Fibonacci indices compatible with a
non-trivial mass spectrum (\(F_1 = F_2 = 1\) gives degenerate masses).
The first non-degenerate branching occurs at depth 3, where the
recursion yields \(\Psi(3) = F_3 = 2\) and \(\Psi(2) = F_2 = 1\). The
prediction is structural: the lightest charged fermion sector uses the
lowest available Fibonacci pair. If a lighter fermionic sector were
discovered, it would need to satisfy \(Q = F_4/(F_4 + F_3) = 3/4\) ---
which is falsifiable.

The same structure extends to quarks, with lower precision:

\begin{longtable}[]{@{}
  >{\raggedright\arraybackslash}p{(\linewidth - 8\tabcolsep) * \real{0.1633}}
  >{\raggedright\arraybackslash}p{(\linewidth - 8\tabcolsep) * \real{0.2653}}
  >{\raggedright\arraybackslash}p{(\linewidth - 8\tabcolsep) * \real{0.2245}}
  >{\raggedright\arraybackslash}p{(\linewidth - 8\tabcolsep) * \real{0.2041}}
  >{\raggedright\arraybackslash}p{(\linewidth - 8\tabcolsep) * \real{0.1429}}@{}}
\toprule\noalign{}
\begin{minipage}[b]{\linewidth}\raggedright
Sector
\end{minipage} & \begin{minipage}[b]{\linewidth}\raggedright
PAC Formula
\end{minipage} & \begin{minipage}[b]{\linewidth}\raggedright
PAC Value
\end{minipage} & \begin{minipage}[b]{\linewidth}\raggedright
Measured
\end{minipage} & \begin{minipage}[b]{\linewidth}\raggedright
Error
\end{minipage} \\
\midrule\noalign{}
\endhead
\bottomrule\noalign{}
\endlastfoot
Charged leptons & \(F_3/(F_3 + F_2) = 2/3\) & 0.666667 & 0.666661 &
\textbf{0.5 ppm} \\
Up-type quarks & \((F_7 - F_3)/F_7 = 11/13\) & 0.846154 & 0.848956 &
\textbf{0.33\%} \\
Down-type quarks & \(\varphi^2/(1+\varphi^2)\) & 0.723607 & 0.731628 &
\textbf{1.10\%} \\
\end{longtable}

The charged-lepton result is essentially exact. The quark results
involve larger Fibonacci indices and transcendental factors, consistent
with the precision hierarchy observed in the gauge couplings.

\textbf{Script}: \texttt{exp\_20\_koide\_formula.py}

\begin{center}\rule{0.5\linewidth}{0.5pt}\end{center}

\subsection{§6. Mass Ratios}\label{mass-ratios}

\subsubsection{§6.1 The Formulas}\label{the-formulas}

Lepton and baryon mass ratios can be expressed as products of Fibonacci
numbers:

\begin{longtable}[]{@{}
  >{\raggedright\arraybackslash}p{(\linewidth - 10\tabcolsep) * \real{0.1061}}
  >{\raggedright\arraybackslash}p{(\linewidth - 10\tabcolsep) * \real{0.1364}}
  >{\raggedright\arraybackslash}p{(\linewidth - 10\tabcolsep) * \real{0.3636}}
  >{\raggedright\arraybackslash}p{(\linewidth - 10\tabcolsep) * \real{0.1364}}
  >{\raggedright\arraybackslash}p{(\linewidth - 10\tabcolsep) * \real{0.1515}}
  >{\raggedright\arraybackslash}p{(\linewidth - 10\tabcolsep) * \real{0.1061}}@{}}
\toprule\noalign{}
\begin{minipage}[b]{\linewidth}\raggedright
Ratio
\end{minipage} & \begin{minipage}[b]{\linewidth}\raggedright
Formula
\end{minipage} & \begin{minipage}[b]{\linewidth}\raggedright
Fibonacci Decomposition
\end{minipage} & \begin{minipage}[b]{\linewidth}\raggedright
Numeric
\end{minipage} & \begin{minipage}[b]{\linewidth}\raggedright
Measured
\end{minipage} & \begin{minipage}[b]{\linewidth}\raggedright
Error
\end{minipage} \\
\midrule\noalign{}
\endhead
\bottomrule\noalign{}
\endlastfoot
\(m_\mu/m_e\) & \(F_4 \times F_6^2 \times (1 + 1/F_7)\) &
\(3 \times 64 \times 14/13\) & 206.769 & 206.768 & \textbf{5 ppm} \\
\(m_\tau/m_e\) & \(F_4 \times F_7 \times F_{11} + F_5\) &
\(3 \times 13 \times 89 + 5\) & 3476 & 3477.23 & \textbf{0.035\%} \\
\(m_p/m_e\) & \(F_4 \times F_9 \times F_{12}/F_6\) &
\(3 \times 34 \times 144/8\) & 1836 & 1836.15 & \textbf{0.0083\%} \\
\end{longtable}

A structural asymmetry should be noted. The gauge coupling formulas are
clean 2-number ratios: \(F_4/F_7 = 3/13\) for \(\sin^2\theta_W\),
\(F_4/(2\varphi F_6)\) for \(\alpha_s\). The mass formulas require
products of 3--4 Fibonacci numbers, with the tau formula including an
additive correction (\(+ F_5\)). This may reflect genuine structural
complexity in the mass sector --- masses involve the full depth of the
hierarchy tree while couplings are ratios at fixed levels.
Alternatively, it may indicate that the correct closed forms for mass
ratios have not yet been found, and simpler expressions exist at deeper
Fibonacci indices.

\subsubsection{§6.2 Structural
Observations}\label{structural-observations}

Three features of these formulas are notable:

\begin{enumerate}
\def\labelenumi{\arabic{enumi}.}
\item
  \textbf{\(F_4 = 3\) appears in every formula.} The muon ratio has it
  as the leading factor. The tau ratio has it as the leading factor. The
  proton ratio has it as the leading factor. If \(F_4\) encodes the
  number of fermion generations (which the Standard Model does not
  explain), its universal presence in mass formulas would follow from
  the branching structure --- each mass ratio involves multiplication
  across the full generation depth.
\item
  \textbf{Cross-consistency.} The derived ratios are consistent:

  \begin{longtable}[]{@{}llll@{}}
  \toprule\noalign{}
  Derived Ratio & From Formulas & Measured & Error \\
  \midrule\noalign{}
  \endhead
  \bottomrule\noalign{}
  \endlastfoot
  \(m_\tau/m_\mu\) & \((m_\tau/m_e)/(m_\mu/m_e)\) & 16.817 &
  \textbf{0.036\%} \\
  \(m_p/m_\mu\) & \((m_p/m_e)/(m_\mu/m_e)\) & 8.880 &
  \textbf{0.009\%} \\
  \end{longtable}
\item
  \textbf{Fibonacci index range.} The formulas use Fibonacci numbers
  from \(F_4\) through \(F_{12}\). Higher mass ratios require higher
  Fibonacci indices, consistent with deeper hierarchy levels encoding
  heavier particles.
\end{enumerate}

\subsubsection{§6.3 The Precision Pattern}\label{the-precision-pattern}

\begin{longtable}[]{@{}lll@{}}
\toprule\noalign{}
Ratio & Precision & Best Fibonacci Numbers Used \\
\midrule\noalign{}
\endhead
\bottomrule\noalign{}
\endlastfoot
\(m_\mu/m_e\) & 5 ppm & \(F_4, F_6, F_7\) \\
\(m_p/m_e\) & 83 ppm & \(F_4, F_6, F_9, F_{12}\) \\
\(m_\tau/m_e\) & 350 ppm & \(F_4, F_5, F_7, F_{11}\) \\
\end{longtable}

The muon-electron ratio achieves the highest precision (matching the
\(\alpha\) result), while the tau formula involves an additive
correction (\(+ F_5\)) that may indicate a more complex structure.

\subsubsection{§6.4 Methodology
Transparency}\label{methodology-transparency}

The mass ratio formulas presented above were found by \textbf{systematic
search} through products and ratios of Fibonacci numbers, then validated
against measured values. They are pattern-matches --- not predictions
derived from first principles. The search template
\(F_a \times F_b^c \times (1 + 1/F_d)\) was tested across Fibonacci
indices \(F_3\) through \(F_{13}\), and the best-matching combinations
are reported. This is an important distinction from the gauge coupling
results, where the formulas have structural motivation (Noether charges,
gauge group dimensions). Until the mass formulas can be derived from the
PAC Lagrangian (§2.1) without searching, they should be treated as
empirical observations compatible with Fibonacci structure, not as
structural predictions.

A stoichiometric analysis (milestone3/exp\_13) found that \(F_4 = 3\) is
\textbf{6,111× more selective} than the next-best integer substitution
across the combined formula set --- a 99.98th-percentile result. This
supports the structural significance of \(F_4\) but does not upgrade
pattern-matching to derivation.

\subsubsection{§6.5 Falsification}\label{falsification}

A Monte Carlo test generated \(10{,}000\) random formulas using the same
structural template (products and ratios of Fibonacci numbers \(F_3\)
through \(F_{13}\)) and checked how many simultaneously match all three
mass ratios within the observed precision.

\textbf{Result}: Zero out of \(10{,}000\) random combinations matched.
Joint probability: \(p < 10^{-4}\).

\textbf{Caveat on sample size}: The formula template
\(F_a \times F_b^c \times (1 + 1/F_d)\) with Fibonacci indices from
\(F_3\) to \(F_{13}\) spans \(\sim 11^4 \approx 14{,}641\) combinations
for a single ratio. A \(10{,}000\)-trial Monte Carlo samples a
comparable fraction of this space. The \(p < 10^{-4}\) bound is
therefore a lower bound on the joint probability, not a precise
estimate. A \(10^6\)-trial Monte Carlo would strengthen the claim. The
current result is sufficient to reject the null hypothesis (random
Fibonacci combinations match by chance) but should not be
over-interpreted as a precise probability.

The recurrence of \(F_7 = 13\) across the mass formulas (in
\(m_\mu/m_e\) as \(1 + 1/F_7\) and in \(m_\tau/m_e\) as a factor) has
independent probability \(p = 0.014\).

\textbf{Scripts}: \texttt{mass\_derivation/exp\_05\_tighten\_mass.py},
\texttt{mass\_derivation/exp\_06\_validate\_tight.py},
\texttt{mass\_derivation/exp\_04\_mass\_falsification.py}

\begin{center}\rule{0.5\linewidth}{0.5pt}\end{center}

\subsection{§7. Mixing Angles}\label{mixing-angles}

\subsubsection{§7.1 Quark Mixing (CKM
Matrix)}\label{quark-mixing-ckm-matrix}

The Cabibbo angle \(\theta_C\) is the dominant off-diagonal element of
the CKM matrix:

\[\theta_C = \arctan\left(\frac{F_4}{F_7}\right) = \arctan\left(\frac{3}{13}\right) = 12.995°\]

\textbf{Measured}: \(\theta_C = 13.00° \pm 0.05°\)\\
\textbf{Error}: \(< 0.05°\)

This is the same ratio \(F_4/F_7 = 3/13\) that gives \(\sin^2\theta_W\).
The Weinberg and Cabibbo angles share arithmetic:
\(\sin^2\theta_W \approx \tan\theta_C\) (0.23121 vs 0.23092,
\(0.4\sigma\) agreement). This relationship is not predicted by the
Standard Model.

\subsubsection{§7.2 Neutrino Mixing (PMNS
Matrix)}\label{neutrino-mixing-pmns-matrix}

Neutrino oscillation angles follow from Fibonacci ratios at different
levels:

\begin{longtable}[]{@{}
  >{\raggedright\arraybackslash}p{(\linewidth - 8\tabcolsep) * \real{0.1373}}
  >{\raggedright\arraybackslash}p{(\linewidth - 8\tabcolsep) * \real{0.2549}}
  >{\raggedright\arraybackslash}p{(\linewidth - 8\tabcolsep) * \real{0.2157}}
  >{\raggedright\arraybackslash}p{(\linewidth - 8\tabcolsep) * \real{0.1961}}
  >{\raggedright\arraybackslash}p{(\linewidth - 8\tabcolsep) * \real{0.1961}}@{}}
\toprule\noalign{}
\begin{minipage}[b]{\linewidth}\raggedright
Angle
\end{minipage} & \begin{minipage}[b]{\linewidth}\raggedright
PAC Formula
\end{minipage} & \begin{minipage}[b]{\linewidth}\raggedright
Prediction
\end{minipage} & \begin{minipage}[b]{\linewidth}\raggedright
Measured
\end{minipage} & \begin{minipage}[b]{\linewidth}\raggedright
\(\Delta\)
\end{minipage} \\
\midrule\noalign{}
\endhead
\bottomrule\noalign{}
\endlastfoot
\(\theta_{12}\) (solar) & \(\arctan(F_3/F_4) = \arctan(2/3)\) & 33.69° &
\(33.41° \pm 0.4°\) & \textbf{0.28°} \\
\(\theta_{13}\) (reactor) & \(\arctan(F_3/F_7) = \arctan(2/13)\) & 8.75°
& \(8.54° \pm 0.2°\) & \textbf{0.21°} \\
\(\theta_{23}\) (atmospheric) & 45° (maximal) & 45° & \(49.0° \pm 1°\) &
4° \\
\end{longtable}

\(\theta_{12}\) and \(\theta_{13}\) are both within \$\sim\(1\)\sigma\$
of the PAC predictions. \(\theta_{23}\) deviates by 4° from maximal
mixing. This is the weakest prediction in this section. Either the
atmospheric angle receives corrections beyond leading-order PAC, or the
prediction is wrong.

\subsubsection{§7.3 Lepton-Quark
Hierarchy}\label{lepton-quark-hierarchy}

The ratio of lepton to quark mixing angles:

\[\frac{\theta_{12}^{\text{PMNS}}}{\theta_{12}^{\text{CKM}}} = \frac{33.41°}{13.00°} = 2.570\]

The PAC prediction: \(\varphi^2 = 2.618\).\\
\textbf{Agreement}: \(0.8\sigma\)

This suggests that leptons and quarks are separated by exactly two
levels in the PAC hierarchy tree. The lepton mixing angle is
\(\varphi^2 \approx 2.6\) times the quark mixing angle because the PMNS
matrix operates two hierarchy levels above the CKM matrix.

\textbf{Scripts}: \texttt{validated/37\_tree\_geometry.py},
\texttt{validated/38\_phi\_squared\_discovery.py}

\begin{center}\rule{0.5\linewidth}{0.5pt}\end{center}

\subsection{§8. Bell Correlations and
Entanglement}\label{bell-correlations-and-entanglement}

\subsubsection{§8.1 The PAC Entanglement
Limit}\label{the-pac-entanglement-limit}

Fibonacci parent-child amplitudes in a two-particle entangled state:

\[a_1 = \frac{1}{\sqrt{1+\varphi^2}}, \quad a_2 = \frac{\varphi}{\sqrt{1+\varphi^2}}\]

The squared visibility:

\[(2a_1 a_2)^2 = \frac{4\varphi^2}{(1+\varphi^2)^2}\]

Using \(1 + \varphi^2 = 2 + \varphi\) (from \(\varphi^2 = \varphi + 1\))
and \((2+\varphi)^2 = 5(1+\varphi) = 5\varphi^2\):

\[(2a_1 a_2)^2 = \frac{4(\varphi + 1)}{5(1 + \varphi)} = \frac{4}{5} \qquad \text{(algebraically exact)}\]

This is not a numerical coincidence. It is an identity of the golden
ratio.

\subsubsection{§8.2 The Bell Parameter}\label{the-bell-parameter}

The PAC Bell parameter from the charged-lepton sector alone:

\[S_{\text{lepton}} = 2\sqrt{1 + \frac{4}{5}} = 2\sqrt{\frac{9}{5}} = \frac{6}{\sqrt{5}} \approx 2.683\]

The quantum mechanical maximum is
\(S_{\text{QM}} = 2\sqrt{2} \approx 2.828\). The experimental value
(Storz et al., 2023 {[}4{]}) is \(S_{\text{exp}} = 2.79 \pm 0.03\).

The gap in squared Bell parameters:

\[(2\sqrt{2})^2 - \left(\frac{6}{\sqrt{5}}\right)^2 = 8 - \frac{36}{5} = \frac{4}{5}\]

\subsubsection{§8.3 The Combined
Prediction}\label{the-combined-prediction}

The charged-lepton contribution to the squared visibility is \(4/5\). If
the neutrino sector contributes the remaining \(1/5\), the total squared
visibility is:

\[V^2_{\text{total}} = \frac{4}{5} + \frac{1}{5} = 1\]

The additivity of \(V^2\) across sectors follows if the charged-lepton
and neutrino channels contribute independently to the entanglement
budget, as expected for orthogonal weak-isospin eigenstates whose
amplitudes do not interfere.

The combined PAC Bell parameter is then:

\[S_{\text{PAC}} = 2\sqrt{1 + V^2_{\text{total}}} = 2\sqrt{1 + 1} = 2\sqrt{2} \approx 2.828\]

This is the quantum mechanical maximum exactly. When both sectors are
included, PAC recovers the Tsirelson bound \(S = 2\sqrt{2}\). The
framework does not predict a deviation from quantum mechanics --- it
partitions the quantum maximum into a charged-lepton fraction (\(4/5\))
and a neutrino fraction (\(1/5\)), both determined algebraically by the
golden ratio.

The experimental value \(S_{\text{exp}} = 2.79 \pm 0.03\) is
\(1.3\sigma\) below \(2\sqrt{2}\), consistent with both the quantum
maximum and with a measurement that predominantly samples the
charged-lepton channel (which would give \(6/\sqrt{5} \approx 2.683\) in
the absence of neutrino contributions).

\textbf{Scripts}: \texttt{validated/32\_pac\_bell\_deep\_dive.py},
\texttt{validated/36\_bell\_neutrino\_resolution.py}

\begin{center}\rule{0.5\linewidth}{0.5pt}\end{center}

\subsection{§9. Turbulence}\label{turbulence}

\subsubsection{§9.1 She-Lévêque
Intermittency}\label{she-luxe9vuxeaque-intermittency}

The She-Lévêque formula describes intermittent fluctuations in turbulent
flows:

\[\zeta_p = \frac{p}{k} + C_0 \left[1 - \left(\frac{C_0}{C_0 + 1}\right)^{p/d}\right]\]

where \(k\), \(C_0\), and \(d\) are parameters traditionally fitted to
experimental data {[}3{]}.

In Fibonacci arithmetic, all three parameters are determined:

\begin{longtable}[]{@{}llll@{}}
\toprule\noalign{}
Parameter & Standard Notation & Fibonacci & Value \\
\midrule\noalign{}
\endhead
\bottomrule\noalign{}
\endlastfoot
\(k\) (scaling) & 9 & \((F_4)^2 = 3^2\) & 9 \\
\(C_0\) (coefficient) & 2 & \(F_3\) & 2 \\
\(d\) (exponent base) & 3 & \(F_4\) & 3 \\
\(\beta\) (cascade ratio) & 2/3 & \(F_3/F_4\) & 0.667 \\
\end{longtable}

The She-Lévêque formula becomes:

\[\zeta_p = \frac{p}{9} + 2\left[1 - \left(\frac{2}{3}\right)^{p/3}\right]\]

Every parameter is a Fibonacci number or ratio. The cascade ratio
\(\beta = 2/3\) is the same \(F_3/F_4\) that gives the Koide formula.

\subsubsection{\texorpdfstring{§9.2 The \(k = d \times F_{d+1}\)
Derivation}{§9.2 The k = d \textbackslash times F\_\{d+1\} Derivation}}\label{the-k-d-times-f_d1-derivation}

Why \(k = 9\)? The dimensional formula:

\[k(d) = d \times F_{d+1}\]

\begin{longtable}[]{@{}llll@{}}
\toprule\noalign{}
Dimension & \(k\) & Calculation & Verification \\
\midrule\noalign{}
\endhead
\bottomrule\noalign{}
\endlastfoot
\(d = 2\) & 4 & \(2 \times F_3 = 2 \times 2\) & 2D enstrophy cascade \\
\(d = 3\) & 9 & \(3 \times F_4 = 3 \times 3\) & She-Lévêque (3D) \\
\(d = 4\) & 20 & \(4 \times F_5 = 4 \times 5\) & \textbf{Prediction} \\
\end{longtable}

The 2D result (\(k = 4\)) was verified independently: the 2D turbulence
spectrum uses \(\beta = F_4/F_5 = 3/5\) (within 2\% of experiment), one
Fibonacci index higher than the 3D cascade, consistent with the
enstrophy cascade requiring an additional scale level.

The 4D prediction (\(k = 20\)) is falsifiable through 4D turbulence
simulations.

\subsubsection{§9.3 The Kolmogorov
Connection}\label{the-kolmogorov-connection}

The Kolmogorov \(-5/3\) exponent for the energy spectrum:

\[E(k) \propto k^{-5/3}\]

The exponent \(5/3\) is derived from dimensional analysis (Kolmogorov,
1941) and is not in dispute. The observation here is that
\(5/3 = F_5/F_4\): it is a ratio of consecutive Fibonacci numbers.
Similarly, the energy dissipation scaling \(\epsilon^{2/3}\) involves
\(2/3 = F_3/F_4\). These are factual restatements --- we note the
Fibonacci structure of known results, not a new derivation of the
Kolmogorov exponent. The question is whether this Fibonacci structure is
coincidental (small-integer ratios are common in dimensional analysis)
or connected to the She-Lévêque results above, where the Fibonacci
structure is less trivially embedded.

\textbf{Scripts}: \texttt{exp\_21\_she\_leveque.py},
\texttt{exp\_11\_k9\_derivation.py}, \texttt{exp\_12\_falsification.py}

\subsubsection{§9.4 Wilson-Fisher Critical
Exponents}\label{wilson-fisher-critical-exponents}

The Wilson-Fisher fixed point governs the 3D Ising universality class
--- the critical behaviour of phase transitions in magnets, fluids, and
binary alloys. The correlation length exponent \(\nu\) determines how
the correlation length diverges near criticality:

\[\xi \sim |T - T_c|^{-\nu}\]

The conformal bootstrap value is \(\nu = 0.6299709 \pm 0.0000040\)
{[}11{]}.

In PAC terms:

\[\nu = \frac{F_3}{F_4 \cdot \Xi} = \frac{2}{3\Xi} = 0.629865\ldots\]

\textbf{Error}: 0.017\%

The formula decomposes as a product of two independently motivated
quantities:

\begin{itemize}
\tightlist
\item
  \textbf{\(2/3 = F_3/F_4\)}: the E-I-S cycle ratio --- the same ratio
  that gives the Koide formula (§5) and the She-Lévêque cascade (§9.1)
\item
  \textbf{\(1/\Xi\)}: the reciprocal of the balance constant (Paper 2)
\end{itemize}

If this decomposition is not coincidental, the Wilson-Fisher exponent
reflects the E-I-S cycle topology at the SEC balance point: correlation
length divergence at a phase transition would be the PAC recursion
running at \(2/3\) of the balance-mediated rate.

\textbf{Broader exponent analysis.} A systematic search across 7
Wilson-Fisher critical exponents found Fibonacci/PAC expressions within
1\% for 6 of the 7:

\begin{longtable}[]{@{}
  >{\raggedright\arraybackslash}p{(\linewidth - 8\tabcolsep) * \real{0.1818}}
  >{\raggedright\arraybackslash}p{(\linewidth - 8\tabcolsep) * \real{0.2000}}
  >{\raggedright\arraybackslash}p{(\linewidth - 8\tabcolsep) * \real{0.2909}}
  >{\raggedright\arraybackslash}p{(\linewidth - 8\tabcolsep) * \real{0.2000}}
  >{\raggedright\arraybackslash}p{(\linewidth - 8\tabcolsep) * \real{0.1273}}@{}}
\toprule\noalign{}
\begin{minipage}[b]{\linewidth}\raggedright
Exponent
\end{minipage} & \begin{minipage}[b]{\linewidth}\raggedright
Literature
\end{minipage} & \begin{minipage}[b]{\linewidth}\raggedright
PAC Expression
\end{minipage} & \begin{minipage}[b]{\linewidth}\raggedright
PAC Value
\end{minipage} & \begin{minipage}[b]{\linewidth}\raggedright
Error
\end{minipage} \\
\midrule\noalign{}
\endhead
\bottomrule\noalign{}
\endlastfoot
\(\nu\) & 0.6300 & \(F_3/(F_4 \cdot \Xi)\) & 0.6299 &
\textbf{0.017\%} \\
\(\eta\) & 0.0363 & \(2/F_{10}\) & 0.0364 & \textbf{0.18\%} \\
\(\gamma\) & 1.2371 & \(2/\varphi\) & 1.2361 & \textbf{0.083\%} \\
\(\delta\) & 4.789 & \(\varphi + \pi\) & 4.760 & \textbf{0.62\%} \\
\(\omega\) & 0.8297 & \(\ln\varphi / \gamma_{\mathrm{E}}\) & 0.8337 &
\textbf{0.48\%} \\
\(\beta\) & 0.3265 & \(\varphi/F_5\) & 0.3236 & \textbf{0.89\%} \\
\(\alpha_{\text{heat}}\) & 0.1096 & --- & --- & no match \\
\end{longtable}

The specific heat exponent \(\alpha_{\text{heat}}\) admits no Fibonacci
expression within 1\%. This is recorded as an honest limitation.

\textbf{Null test.} A Monte Carlo search of 23,376 formula candidates
across 2,000 random-constant trials found a mean of 1.89 hits within
1\%, compared to 20 hits for the PAC constants (\(p = 0.0000\)).

The best alternative to \(\nu = 2/(3\Xi)\) is \(\pi/(3\Xi)\), which
matches at 1.06\% --- 63× worse. The factor \(2/3\) is uniquely
selected.

\textbf{Script}: \texttt{milestone3/scripts/exp\_07\_wilson\_fisher.py}

\begin{center}\rule{0.5\linewidth}{0.5pt}\end{center}

\subsection{§10. Casimir Effect}\label{casimir-effect}

\subsubsection{§10.1 The Factor 240}\label{the-factor-240}

The Casimir force between two parallel conducting plates:

\[\frac{F}{A} = -\frac{\pi^2 \hbar c}{240\,d^4}\]

The factor 240 arises from zeta function regularisation:
\(\sum n^3 \to \zeta(-3) = 1/120\), combined with two-plate symmetry
giving \(\pi^2/240\).

In Fibonacci arithmetic:

\[240 = F_3 \times F_4 \times F_5 \times F_6 = 2 \times 3 \times 5 \times 8\]

Four consecutive Fibonacci numbers. The sub-factor
\(120 = F_4 \times F_5 \times F_6 = 3 \times 5 \times 8\) (three
consecutive) gives \(\zeta(-3)\); the Casimir force gradient adds
\(F_3 = 2\) through differentiation.

\subsubsection{§10.2 Mersenne Dimensions}\label{mersenne-dimensions}

The zeta function regularisation denominators show Fibonacci product
structure only at Mersenne dimensions \(d = 2^n - 1\):

\begin{longtable}[]{@{}llll@{}}
\toprule\noalign{}
\(d\) & \(\zeta(-d)\) denominator & Fibonacci product? & Mersenne? \\
\midrule\noalign{}
\endhead
\bottomrule\noalign{}
\endlastfoot
1 & 12 & \(F_3^2 \times F_4\) & \(2^1 - 1\) ✓ \\
3 & 120 & \(F_4 \times F_5 \times F_6\) & \(2^2 - 1\) ✓ \\
5 & 252 & No (factor 7) & ✗ \\
7 & 240 & \(F_3 \times F_4 \times F_5 \times F_6\) & \(2^3 - 1\) ✓ \\
9 & 132 & No (factor 11) & ✗ \\
\end{longtable}

The pattern: vacuum regularisation has Fibonacci structure at the exact
dimensions relevant to string theory (\(d = 1\)), the physical Casimir
effect (\(d = 3\)), and M-theory's extra dimensions (\(d = 7\)).
Non-Mersenne dimensions have non-Fibonacci prime factors.

\textbf{Scripts}: \texttt{exp\_15\_casimir\_sec.py},
\texttt{exp\_16\_mersenne\_verification.py}

\begin{center}\rule{0.5\linewidth}{0.5pt}\end{center}

\subsection{§11. Gravity Hierarchy}\label{gravity-hierarchy}

The ratio of electromagnetic to gravitational force between two protons:

\[\frac{F_{\text{EM}}}{F_{\text{grav}}} = \frac{e^2/(4\pi\epsilon_0)}{Gm_p^2} \approx 1.24 \times 10^{38}\]

The Fibonacci prediction uses the gauge-squared depth:

\[183 = F_7^2 + F_7 + 1 = 169 + 13 + 1\]

Then \(F_{183} \approx 1.27 \times 10^{38}\).

This is an order-of-magnitude result, not a precision prediction. The
hierarchy ratio is \$\sim\$10³⁸, and \(F_{183}\) is \$\sim\$10³⁸. The
match is suggestive but not compelling by itself.

What makes it non-trivial is the structural argument: if the EM
hierarchy depth is \(F_{10} = 55\) (from the \(\alpha\) formula), and
gravity requires an additional gauge-squared depth
\(F_7^2 + F_7 + 1 = 183\), then the gravitational hierarchy is encoded
by \(F_{183}\) --- the Fibonacci number at a depth determined by the
gauge closure constant \(F_7 = 13\).

\textbf{On uniqueness}: The Fibonacci sequence grows as
\(F_n \approx \varphi^n/\sqrt{5}\), so
\(\log_{10}(F_n) \approx 0.2090n\). At \(n = 183\), this gives
\(\sim 10^{38.2}\). Nearby indices (\(n = 178 \to 10^{37.2}\),
\(n = 188 \to 10^{39.3}\)) are within an order of magnitude. The
falsification test is therefore not about whether \(F_{183}\) is
uniquely close to \(10^{38}\) --- several nearby Fibonacci numbers are
comparably close. The test is whether the structural formula
\(n = F_7^2 + F_7 + 1 = 183\) is uniquely motivated. The polynomial
\(N^2 + N + 1\) counts elements in the projective plane \(PG(2, N)\);
evaluated at the gauge closure constant \(N = F_7 = 13\), it gives
precisely the depth that encodes the EM-gravity hierarchy. No other
simple polynomial in \(F_7\) produces a depth whose Fibonacci number
matches the hierarchy ratio.

\textbf{Milestone3 update (exp\_23, exp\_26).} The precise gap between
\(\log_{10}(F_{183})\) and \(\log_{10}((M_{\mathrm{Pl}}/m_p)^2)\) is
0.333 in \(\log_{10}\) (factor \$\sim\$2.155). The best correction term
is \(1 + F_{13}/(\pi \cdot F_6^2)\), which brings the residual to
\(7.8 \times 10^{-4}\) in \(\log_{10}\). Among Fibonacci depths
\(F_k^2 + F_k + 1\), only \(k = 7\) (depth 183) falls within 0.5 of the
target: rank \#1 cyclotomic depth, with a 40× gap to the next closest. A
Monte Carlo test of 5,000 random integer sequences found
\textbf{0/5,000} matching both \(\alpha_{\mathrm{EM}}\) and gravity
simultaneously using the correction template \(F_a/(m\pi F_b^2)\),
suggesting the joint constraint is structurally non-trivial.

\textbf{Scripts}: \texttt{exp\_23\_gravity\_depth.py},
\texttt{exp\_24\_hierarchy\_f183.py},
\texttt{exp\_26\_hierarchy\_falsification.py}

\begin{center}\rule{0.5\linewidth}{0.5pt}\end{center}

\subsection{§12. What This Is Not}\label{what-this-is-not}

\subsubsection{§12.1 It Is Not Numerology}\label{it-is-not-numerology}

The standard objection to expressing physical constants in terms of
small integers is that small integers are common, and matches are
therefore unsurprising. This objection is valid for individual matches.

Consider: the probability that a randomly chosen integer between 1 and
100 is a Fibonacci number is approximately 16\% (12 Fibonacci numbers
below 100). A single match --- say, the number of \(SU(3)\) generators
being \(F_6 = 8\) --- is not statistically meaningful.

The claim here is not about individual matches. It is about joint
constraints. A single recursion \(\Psi(k) = \Psi(k+1) + \Psi(k+2)\)
selects the Fibonacci sequence. From that sequence:

\begin{itemize}
\tightlist
\item
  \(F_7 = 13 = 1 + 3 + 8 + 1\) (gauge closure) \emph{selects the gauge
  group structure}
\item
  \(F_4/F_7 = 3/13\) (weak mixing) \emph{predicts a coupling constant
  ratio to 0.19\%}
\item
  \(F_3/(F_4 \cdot \varphi \cdot F_{10}) \times \text{correction}\)
  \emph{predicts \(\alpha\) to 5.7 ppm}
\item
  \(F_3/(F_3 + F_2) = 2/3\) \emph{matches the Koide formula to 0.5 ppm}
\item
  \(F_4 \times F_6^2 \times (1 + 1/F_7)\) \emph{matches the
  muon-electron mass ratio to 5 ppm}
\end{itemize}

These are not independent fits. They use the same Fibonacci numbers
(\(F_3, F_4, F_6, F_7, F_{10}\)) in interlocking formulas that must be
simultaneously consistent. The joint probability against chance,
estimated by Monte Carlo (\(10{,}000\) random formula sets with the same
template), is \(p < 10^{-5}\).

\textbf{Cross-domain independence audit (milestone3/exp\_10).} The
results in this paper span particle physics (§§4--7), quantum mechanics
(§8), fluid dynamics (§9.1--9.3), statistical mechanics (§9.4), and
quantum field theory (§10). A structural correlation test identifies
which claims share Fibonacci indices and which are genuinely
independent.

Of 14 claims analysed across the PACSeries, the effective number of
independent degrees of freedom is 7.9 (independence ratio: 0.56). The
top correlated pair is \(\alpha\) and \(\sin^2\theta_W\), which share
\(F_4\) and \(F_7\) --- these are \emph{not} independent results.
Conversely, the turbulence parameters (§9.1--9.3), the Casimir factor
(§10), and the Wilson-Fisher exponent (§9.4) use Fibonacci numbers and
constants (\(\Xi\), \(\gamma_{\mathrm{E}}\), \(\ln\varphi\)) that do not
appear in any particle physics formula.

Grouping claims into 5 conservative independence classes and combining
one \(p\)-value per group via Fisher's method gives a joint significance
of \(p \approx 10^{-147}\), \textbf{conditional on this analysis
structure} --- i.e., given the specific formula templates, the domain
groupings, and the measured constants tested. This figure does not
account for template selection bias (the templates were chosen because
they work; see §12.5 for the α look-elsewhere analysis that quantifies
this for one case). The correction from the naive (uncorrelated)
estimate to the conservative (grouped) estimate is \$\sim\$48 orders of
magnitude --- showing that the analysis properly penalises shared
structure. Even after this penalty and the template-selection caveat,
the joint significance substantially exceeds conventional thresholds. A
global template richness audit (milestone3/exp\_32) tested 50
dimensionless PDG constants --- 7 claimed and 43 unclaimed --- against
the same combined 2-, 3-, and 4-factor Fibonacci template pool
(\textasciitilde26,700 unique values). At 1\% precision, 91\% of
unclaimed constants can be matched, confirming that template-level
matches at this threshold are cheap. At 100 ppm, only 19\% of unclaimed
constants are reachable, versus 43\% of claimed constants. The
claimed/unclaimed median-error ratio is 1.9×. \textbf{Conclusion}:
individual matches at 1\% carry little weight; the joint constraint and
sub-100 ppm results (§15.1--15.2) are where the signal resides.

\textbf{Script}:
\texttt{milestone3/scripts/exp\_10\_independence\_audit.py}

\subsubsection{§12.2 It Is Not a
Derivation}\label{it-is-not-a-derivation}

We do not derive the Standard Model from PAC. A derivation would produce
the gauge group structure, the coupling constants, and the mass spectrum
from axioms alone, with no appeal to measured values. What we present is
weaker: given the PAC recursion, we show that measured values are
consistent with Fibonacci expressions, with precisions that range from
0.5 ppm to 1.7\%.

The gap between ``consistent with'' and ``derived from'' is important. A
derivation of \(\sin^2\theta_W = 3/13\) would explain why the
electroweak mixing occurs at this ratio. Our result does not --- it
observes the coincidence and reports it. The structural arguments (gauge
group selection, Noether charges, hierarchy depths) are suggestive but
do not constitute a proof.

\subsubsection{§12.4 Honest Failures}\label{honest-failures}

Three milestone3 experiments produced results that limit the framework's
claims:

\begin{enumerate}
\def\labelenumi{\arabic{enumi}.}
\item
  \textbf{Null-space prediction failure (exp\_16).} An attempt to use
  the PAC framework to \emph{predict} unmeasured ratios from the null
  space of the formula matrix scored 0/4. Fibonacci-indexed predictions
  showed no enrichment over random integer-indexed formulas
  (\(z = 0.455\), \(p_{\text{MC}} = 0.794\)). The framework
  \textbf{describes} measured constants well but does \textbf{not
  predict} unmeasured ones with current methods.
\item
  \textbf{Crystallisation order is basis-independent (exp\_19,
  FALSIFIED).} The order in which formula constants converge to
  measurement during a PAC-guided optimisation is \textbf{identical}
  across Fibonacci, Lucas, prime, tribonacci, and random sequences.
  There is no Fibonacci-specific dynamics in the convergence process;
  the crystallisation order is a property of the formula structure
  itself.
\item
  \textbf{PAC-Lazy signal is fragile (exp\_24).} The KL-divergence
  discrimination between PAC-matched and unmatched formulas
  (\(p = 0.035\), exp\_21) does not survive bootstrap analysis: the 95\%
  CI on the KL difference is \([-0.044, +0.013]\), \textbf{crossing
  zero}. The signal concentrates in a single formula and does not reduce
  the effective dimensionality below null-space degrees of freedom. This
  is an engineering observation, not a theoretical claim.
\end{enumerate}

These failures are recorded because they constrain interpretation. The
Fibonacci expressions match measured constants with statistical
significance (\(p < 10^{-5}\) jointly), but the framework does not yet
have predictive power for unmeasured observables, and the Fibonacci
sequence is not uniquely selected by the convergence dynamics.

\subsubsection{§12.5 Look-Elsewhere Analysis for
α}\label{look-elsewhere-analysis-for-ux3b1}

The fine structure constant formula (§4.1) achieves 5.7 ppm precision. A
detailed look-elsewhere analysis (milestone3/exp\_09) tested two
distinct questions about this match:

\textbf{Question 1 (Template-class richness):} Can the formula template
\(k/(m \cdot T \cdot F_i) \times (1 - F_j/(n \cdot U \cdot F_p^q))\)
produce matches to \(\alpha\) by chance? An exhaustive enumeration of
1,640,599 valid formulas found 2 matches within 6 ppm. Since
\$\sim\$1.44 matches are expected at random (binomial), \(p = 0.42\).
\textbf{The template class is rich enough that individual matches are
unremarkable.}

\textbf{Question 2 (Fibonacci index specificity):} Within the
\emph{specific} skeleton of the published formula (with \(\varphi\) and
\(4\pi\) fixed), are the Fibonacci index choices special? A Monte Carlo
test drew 7,272 random index sets from \(F_1\)--\(F_{15}\) using the
exact skeleton. \textbf{Zero matched α.} The specific Fibonacci indices
(\(F_7, F_{10}\)) are fine-tuned within that skeleton.

These are not contradictory --- they answer different questions. The
first shows the template as a whole has enough combinatorial capacity to
hit \(\alpha\)-scale values by chance. The second shows the specific
Fibonacci indices are not random choices within that template. The α
formula is interesting \emph{not} because it is a rare hit from a rich
template, but because the same Fibonacci indices (\(F_4, F_7, F_{10}\))
appear across multiple formulas (Weinberg angle, Koide, lepton masses)
with interlocking constraints.

The correct framing is therefore \textbf{joint constraint significance}
(§12.1), not single-formula significance. The \(p < 10^{-5}\) joint
Monte Carlo is the relevant statistic. The individual \(\alpha\) match,
taken in isolation, does not survive look-elsewhere correction.

\textbf{Bonus finding:} A second formula
\(1/(4\varphi F_8) \times (1 - F_9/(3\pi F_8^2))\) matches \(\alpha\) at
1.09 ppm --- better than the published 5.7 ppm result. This further
demonstrates that individual matches are not unique, reinforcing that
the significance lies in the joint system, not in any single formula.

\subsubsection{§12.3 What Would Make It
Stronger}\label{what-would-make-it-stronger}

The following would elevate these results from ``suggestive pattern'' to
``structural prediction'':

\begin{enumerate}
\def\labelenumi{\arabic{enumi}.}
\item
  \textbf{An independent prediction that is confirmed.} The Z' boson at
  395 GeV (§14) is such a prediction. If found, the Fibonacci structure
  gains credibility. If not found, it does not falsify the coupling
  constant results (the Z' prediction depends on additional
  assumptions), but it weakens the framework.
\item
  \textbf{Extension to all mass ratios.} Currently, three mass ratios
  are well matched. A systematic extension to quark masses, neutrino
  mass squared differences, and hadronic mass ratios would either
  strengthen or falsify the pattern.
\item
  \textbf{Running coupling confirmation.} The running of
  \(\sin^2\theta_W\) from \(M_Z\) to lower energies now confirms
  (milestone3/exp\_08) that \(3/13\) is achieved at \(Q = 82.78\) GeV,
  within 3\% of \(M_W\). The remaining test is whether future precision
  measurements at or near the \(M_W\) scale converge to \(3/13\) rather
  than a nearby value. The MOLLER experiment at JLab and future
  \(e^+e^-\) colliders running at the W threshold would sharpen this
  constraint.
\end{enumerate}

\begin{center}\rule{0.5\linewidth}{0.5pt}\end{center}

\subsection{§13. Falsification
Conditions}\label{falsification-conditions}

These results would be falsified by:

\begin{enumerate}
\def\labelenumi{\arabic{enumi}.}
\item
  \textbf{Alternative Fibonacci combinations matching better.} If a
  different set of Fibonacci numbers reproduces the same constants with
  comparable or better precision, the specific formulas presented here
  lose their uniqueness claim. Current exhaustive search over all
  \((F_m, F_n)\) pairs for \(\alpha\) finds no alternative within
  \(2{,}870\times\) of the best match.
\item
  \textbf{Z' non-observation at HL-LHC.} If no resonance is found near
  395 GeV with the predicted properties by 2030, the gauge closure
  prediction weakens, though it does not falsify the coupling constant
  results directly.
\item
  \textbf{Failure to extend to quarks.} If quark mass ratios cannot be
  expressed in Fibonacci arithmetic at comparable precision to the
  lepton results, the pattern may be limited to the lepton sector and
  therefore less fundamental.
\item
  \textbf{An alternative framework reproducing the same constants.} If a
  non-Fibonacci algebraic structure matches the Standard Model
  parameters at equal or better precision, the claim that Fibonacci
  numbers are structurally significant would be undermined.
\item
  \textbf{Precision measurement refuting \(\sin^2\theta_W = 3/13\).} If
  high-precision measurements of \(\sin^2\theta_W\) at all energy scales
  exclude \(3/13\) as either a tree-level value or a value at any
  specific scale, the result is falsified.
\end{enumerate}

\begin{center}\rule{0.5\linewidth}{0.5pt}\end{center}

\subsection{§14. Predictions}\label{predictions}

\subsubsection{§14.1 Z' Boson}\label{z-boson}

The PAC closure at \(F_7 = 13\) with \(1 + 3 + 8 + 1 = 13\) admits one
interpretation: the ``+1'' is the Higgs. An alternative: the ``+1'' is a
13th gauge boson --- a Z' in a remnant \(U(1)\) from a PAC-compatible
extension.

\begin{longtable}[]{@{}ll@{}}
\toprule\noalign{}
Property & PAC Value \\
\midrule\noalign{}
\endhead
\bottomrule\noalign{}
\endlastfoot
Mass & \(395 \pm 20\) GeV \\
Coupling & \(g_{Z'}/g_Z = 1/13\) \\
Width & \$\sim\$64 MeV (narrow) \\
Cross section & \(1/169\) of standard Z' \\
\end{longtable}

The coupling \(1/13 = 1/F_7\) follows from the hierarchy. The mass
follows from the PAC saturation depth
\(N^* = 3F_{10}/(2\pi) \approx 26\) and the \(SU(2)\) symmetry breaking
scale.

Run 3 data (\$\sim\(300 fb\)\^{}\{-1\}\$ at \(\sqrt{s} = 13.6\) TeV,
completing \$\sim\$2026) may already constrain this parameter space if
reanalysed for narrow, weakly-coupled resonances. The full test requires
HL-LHC luminosities (\(3{,}000\) fb\(^{-1}\), physics runs from
\$\sim\$2030 after Long Shutdown 3). Standard LHC Z' searches exclude
sequential Standard Model Z' bosons below
\$\sim\(5 TeV (ATLAS dilepton: ATLAS-CONF-2019-030 [9]; CMS dimuon: CMS-PAS-EXO-19-019 [10]). However, these exclusions assume standard couplings (\)g\_\{Z'\}
\sim g\_Z\$). At \(g_{Z'}/g_Z = 1/13\), the production cross section is
\(1/169\) of the benchmark, placing a 395 GeV Z' well below current
sensitivity thresholds. The narrow width (\$\sim\$64 MeV, compared to
\(\Gamma_Z \approx 2.5\) GeV) further reduces detection efficiency in
searches optimised for broad resonances. A dedicated low-mass,
narrow-width dilepton search would be required.

\textbf{Script}: \texttt{exp\_34\_zprime\_prediction.py}

\subsubsection{§14.2 Neutrino Mixing
Angles}\label{neutrino-mixing-angles}

\begin{longtable}[]{@{}llll@{}}
\toprule\noalign{}
Angle & PAC Prediction & Current Best & Test \\
\midrule\noalign{}
\endhead
\bottomrule\noalign{}
\endlastfoot
\(\theta_{12}\) & 33.69° & \(33.41° \pm 0.4°\) & JUNO (2025+) \\
\(\theta_{13}\) & 8.75° & \(8.54° \pm 0.2°\) & DUNE (2029+) \\
\end{longtable}

\subsubsection{§14.3 4D Turbulence}\label{d-turbulence}

\[k(4) = 4 \times F_5 = 4 \times 5 = 20\]

Testable through 4D numerical turbulence simulations.

\subsubsection{§14.4 GUT-Scale Coupling
Unification}\label{gut-scale-coupling-unification}

If all three couplings derive from the same PAC recursion, they should
unify at a characteristic PAC scale where
\(\alpha^* = 1/\varphi^3 \approx 0.236\).

\subsubsection{§14.5 Dark Matter Density
(Speculative)}\label{dark-matter-density-speculative}

The formula \(\Omega_c = F_7 \cdot \Xi^2 / F_{10}\) gives a cold dark
matter density of \textbf{0.2648}, compared to the Planck 2018
measurement \(\Omega_c = 0.265 \pm 0.007\). This is 0.079\% from the
central value. An alternative formula \(F_3 \cdot \Xi / F_6\) gives
0.2646 (0.148\% error). Among 590 candidate Fibonacci formulas tested,
only 2 (0.34\%) fall within 0.15\% of the measured value. These results
are speculative: the theoretical motivation for connecting gauge closure
(\(F_7\)) and the balance constant (\(\Xi\)) to cosmological densities
remains to be established.

\begin{center}\rule{0.5\linewidth}{0.5pt}\end{center}

\subsection{§15. Summary of Results}\label{summary-of-results}

Results are grouped by evidential weight. \emph{Structural} results
follow from the PAC recursion without reference to measured values.
\emph{Formula search} results match measured constants through
Fibonacci-template searches and are validated by joint constraints
(§12.1) and null tests, not by individual significance.

\subsubsection{§15.1 Structural Results (derived, not
searched)}\label{structural-results-derived-not-searched}

These follow from the PAC recursion and gauge group arithmetic without
fitting to data:

\begin{longtable}[]{@{}
  >{\raggedright\arraybackslash}p{(\linewidth - 4\tabcolsep) * \real{0.3333}}
  >{\raggedright\arraybackslash}p{(\linewidth - 4\tabcolsep) * \real{0.3333}}
  >{\raggedright\arraybackslash}p{(\linewidth - 4\tabcolsep) * \real{0.3333}}@{}}
\toprule\noalign{}
\begin{minipage}[b]{\linewidth}\raggedright
Result
\end{minipage} & \begin{minipage}[b]{\linewidth}\raggedright
Status
\end{minipage} & \begin{minipage}[b]{\linewidth}\raggedright
Domain
\end{minipage} \\
\midrule\noalign{}
\endhead
\bottomrule\noalign{}
\endlastfoot
\(F_7 = 13 = 1 + 3 + 8 + 1\) & Exact (algebraic) & Gauge closure \\
Fibonacci filter: only SU(2), SU(3) & Exact (algebraic) & Gauge
selection \\
\((2a_1 a_2)^2 = 4/5\) & Exact (algebraic) & Entanglement \\
\(k = d \times F_{d+1}\) & Formula with 4D prediction & Turbulence \\
Koide \(Q = 2/3\) to 0.5 ppm & Exact identity at hierarchy depth 3 &
Mass relation \\
\end{longtable}

\subsubsection{§15.2 High-Precision Formula
Matches}\label{high-precision-formula-matches}

These are matches between Fibonacci expressions and measured constants.
Individually, each could be coincidence (§12.5); their significance
comes from the joint constraint (§12.1):

\begin{longtable}[]{@{}lll@{}}
\toprule\noalign{}
Result & Precision & Domain \\
\midrule\noalign{}
\endhead
\bottomrule\noalign{}
\endlastfoot
\(\alpha\) to 5.7 ppm & Measurement & Gauge coupling \\
\(m_\mu/m_e\) to 5 ppm & Measurement (template search) & Mass ratio \\
\(\nu = 2/(3\Xi)\) & 0.017\% & Wilson-Fisher (critical) \\
\(m_p/m_e\) to 83 ppm & Measurement (template search) & Mass ratio \\
\(\sin^2\theta_W = 3/13\) at \(Q \approx M_W\) & 0.19\% & Gauge
coupling \\
\(\alpha_s\) to 1.71\% & Measurement & Gauge coupling \\
\(m_\tau/m_e\) to 350 ppm & Measurement (template search) & Mass
ratio \\
\(\theta_C\) to \(< 0.05°\) & Measurement & CKM mixing \\
\(\theta_{12}^{\text{PMNS}}\) to 0.28° & Measurement & Neutrino
mixing \\
\(\theta_{13}^{\text{PMNS}}\) to 0.21° & Measurement & Neutrino
mixing \\
\end{longtable}

\subsubsection{§15.3 Small-Integer Compatible (noted, not claimed as
evidence)}\label{small-integer-compatible-noted-not-claimed-as-evidence}

These coincide with Fibonacci numbers but involve small integers where
the match probability is high (\$\sim\$16\%). They are included for
completeness but carry negligible individual weight:

\begin{longtable}[]{@{}
  >{\raggedright\arraybackslash}p{(\linewidth - 4\tabcolsep) * \real{0.3636}}
  >{\raggedright\arraybackslash}p{(\linewidth - 4\tabcolsep) * \real{0.2727}}
  >{\raggedright\arraybackslash}p{(\linewidth - 4\tabcolsep) * \real{0.3636}}@{}}
\toprule\noalign{}
\begin{minipage}[b]{\linewidth}\raggedright
Result
\end{minipage} & \begin{minipage}[b]{\linewidth}\raggedright
Note
\end{minipage} & \begin{minipage}[b]{\linewidth}\raggedright
Domain
\end{minipage} \\
\midrule\noalign{}
\endhead
\bottomrule\noalign{}
\endlastfoot
\(5/3 = F_5/F_4\) & Kolmogorov exponent from dimensional analysis &
Turbulence \\
\(\beta = F_3/F_4 = 2/3\) & She-Lévêque cascade ratio (but see
\(k = d \times F_{d+1}\) above) & Turbulence \\
Casimir 240 = \(F_3 F_4 F_5 F_6\) & Product of consecutive Fibonacci
numbers; Mersenne-dimension restriction (§10.2) is the non-trivial part
& QFT regularisation \\
\end{longtable}

\subsubsection{§15.4 Forward Predictions}\label{forward-predictions}

\begin{longtable}[]{@{}lll@{}}
\toprule\noalign{}
Prediction & Value & Experiment \\
\midrule\noalign{}
\endhead
\bottomrule\noalign{}
\endlastfoot
Z' boson & 395 ± 20 GeV, coupling 1/13 & HL-LHC (\textasciitilde2030) \\
\(k(4) = 20\) & 4D turbulence scaling & Numerical simulation (now) \\
\(\sin^2\theta_W = 3/13\) near \(M_W\) & Q = 82.78 GeV & MOLLER /
FCC-ee \\
Neutrino \(\theta_{12}\) = 33.69° & JUNO & 2025+ \\
Neutrino \(\theta_{13}\) = 8.75° & DUNE & 2029+ \\
\end{longtable}

\begin{center}\rule{0.5\linewidth}{0.5pt}\end{center}

\subsection{§16. Connections to the
PACSeries}\label{connections-to-the-pacseries}

\begin{longtable}[]{@{}
  >{\raggedright\arraybackslash}p{(\linewidth - 2\tabcolsep) * \real{0.3889}}
  >{\raggedright\arraybackslash}p{(\linewidth - 2\tabcolsep) * \real{0.6111}}@{}}
\toprule\noalign{}
\begin{minipage}[b]{\linewidth}\raggedright
Paper
\end{minipage} & \begin{minipage}[b]{\linewidth}\raggedright
Connection
\end{minipage} \\
\midrule\noalign{}
\endhead
\bottomrule\noalign{}
\endlastfoot
Paper 1: Structure Cost of Erasure & \(\alpha\) interpreted as Landauer
payment rate through \(F_{10} = 55\) hierarchy levels; gauge hierarchy
\(\xi(SU(3)) > \xi(SU(2)) > \xi(U(1))\) confirmed at \(p < 10^{-11}\) \\
Paper 2: Balance Constant & Same \(F_{10} = 55\) in
\(\Xi = 1 + \pi/55\); Ξ appears at boundaries where gauge couplings are
evaluated \\
Paper 3: Feigenbaum Constants & Same \(F_{10} = 55\) in closed-form
expressions for \(r_\infty\), \(\delta\), \(|\alpha|\) --- Fibonacci
universality extends from nonlinear dynamics to particle physics \\
\textbf{Paper 4 (this paper)} & \textbf{Standard Model parameters from
Fibonacci arithmetic} \\
Paper 5: Classical Physics & Maxwell's equations as depth-2 PAC
recursion; MED → D = 3; \(k = d \times F_{d+1}\) connects EM and
turbulence \\
Paper 6: Computational Validation & PAC conservation observed in ML
systems validates the operational principle \\
\end{longtable}

The common thread across Papers 1--6 is \(F_7 = 13\) and
\(F_{10} = 55\). These are not free parameters. \(F_7\) is determined by
gauge group arithmetic (\(1 + 3 + 8 + 1 = 13\)). \(F_{10}\) is
determined by the Feigenbaum structure (Paper 3) and the Landauer
hierarchy (Paper 1). Both are structural consequences of the PAC
recursion.

\begin{center}\rule{0.5\linewidth}{0.5pt}\end{center}

\subsection{Acknowledgments}\label{acknowledgments}

\emph{(To be added.)}

\begin{center}\rule{0.5\linewidth}{0.5pt}\end{center}

\subsection{References}\label{references}

\begin{enumerate}
\def\labelenumi{\arabic{enumi}.}
\tightlist
\item
  Particle Data Group (2022). ``Review of Particle Physics.''
  \emph{Prog. Theor. Exp. Phys.}, 2022(8), 083C01.
\item
  Koide, Y. (1982). ``New Formula for the Cabibbo Angle.'' \emph{Phys.
  Rev.~Lett.}, 49, 723--724.
\item
  She, Z.-S. and Lévêque, E. (1994). ``Universal Scaling Laws in Fully
  Developed Turbulence.'' \emph{Phys. Rev.~Lett.}, 72(3), 336--339.
\item
  Storz, S. et al.~(2023). ``Loophole-free Bell inequality violation
  with superconducting circuits.'' \emph{Nature}, 617, 265--270.
\item
  Tiesinga, E. et al.~(2024). ``CODATA recommended values of the
  fundamental physical constants: 2022.'' \emph{J. Phys. Chem. Ref.
  Data}, 53(3), 030801.
\item
  Groom, P. (2026). ``The Structure Cost of Erasure.'' PACSeries Paper
  1. Dawn Field Institute.
\item
  Groom, P. (2026). ``The Balance Constant and Its Decomposition.''
  PACSeries Paper 2. Dawn Field Institute.
\item
  Groom, P. (2026). ``Feigenbaum Constants from Fibonacci Arithmetic.''
  PACSeries Paper 3. Dawn Field Institute.
\item
  ATLAS Collaboration (2019). ``Search for high-mass dilepton resonances
  using 139 fb⁻¹ of pp collision data collected at √s = 13 TeV with the
  ATLAS detector.'' ATLAS-CONF-2019-030.
\item
  CMS Collaboration (2019). ``Search for high-mass resonances in
  dilepton final states in proton-proton collisions at √s = 13 TeV.''
  CMS-PAS-EXO-19-019.
\item
  Kos, F., Poland, D., Simmons-Duffin, D. and Vichi, A. (2016).
  ``Precision islands in the Ising and O(N) models.'' \emph{JHEP},
  2016(8), 36.
\end{enumerate}

\begin{center}\rule{0.5\linewidth}{0.5pt}\end{center}

\emph{All code, data, and experiment scripts for this paper and the full
PACSeries are publicly available at
\url{https://github.com/dawnfield-institute/dawn-field-theory}. See the
accompanying publication package README.md for reproduction
instructions.}

\end{document}
